\section{Análisis Lexicográfico (Pregunta 3)}

\textit{Al aplicar la estrategia lexicográfica directa minimizando sólo $f_2$,
¿qué sucedió con el valor de $f_1$? Compara este valor con el extremo del
frente de Pareto de NSGA-II.}

\subsection{Qué le ocurrió a $f_1$}

Al optimizar \emph{exclusivamente} el Tiempo ($f_2$), el GA llevó la solución
a $\mathbf{x} = (\num{\AprioriLexVarUno}, \num{\AprioriLexVarDos})$, donde
$f_2$ alcanza su mínimo global $\num{\AprioriLexTiempo}$. Pero al ignorar por
completo al Costo, éste se disparó hasta su \textbf{valor máximo posible} en
el dominio:
\[
  f_1 = \num{\AprioriLexCosto}, \qquad f_2 = \num{\AprioriLexTiempo}.
\]
Esto es exactamente lo que predice la teoría: minimizar $f_2 = (x_1-5)^2 +
(x_2-5)^2$ empuja la solución a la esquina $(5,5)$, que es justamente el punto
donde $f_1 = 4(25)+4(25) = \num{200}$ es \emph{máximo}. La prioridad rígida
sobre un objetivo sacrifica por completo al otro.

\subsection{Comparación con el extremo del frente NSGA-II}

El frente aproximado por NSGA-II tiene su extremo de menor tiempo en
$f_1 = \num{\MetMaxFuno}$, $f_2 = \num{\MetMinFdos}$ (la solución
\texttt{post\_lexico} de la Tabla~\ref{tab:soluciones}). Como se aprecia en la
Figura~\ref{fig:zoom}, \textbf{ambos enfoques lexicográficos aterrizan en el
mismo punto} $(\num{200}, \num{0})$: la decisión lexicográfica a priori y la
selección lexicográfica a posteriori son indistinguibles aquí, porque ambas
persiguen el mismo objetivo extremo del frente.

\begin{figure}[h]
  \centering
  \includegraphics[width=0.85\linewidth]{pareto_front_zoom.png}
  \caption{Zoom al codo del frente. Las soluciones lexicográficas (cuadrado
  amarillo y marcador negro) coinciden en el extremo $(\num{200}, \num{0})$,
  el punto de mayor costo del frente.}
  \label{fig:zoom}
\end{figure}

\subsection{El costo de la prioridad rígida}

La lección es contundente: pasar de la solución \emph{ponderada}
($f_1 \approx \num{\AprioriPondCosto}$) a la \emph{lexicográfica}
($f_1 = \num{\AprioriLexCosto}$) implica multiplicar el costo operativo por un
factor superior a $\times 10$ a cambio de recortar el tiempo desde
$\approx \num{\AprioriPondTiempo}$ hasta $\num{\AprioriLexTiempo}$. El enfoque
lexicográfico es apropiado únicamente cuando el objetivo prioritario es
verdaderamente \textbf{innegociable} (p.\ ej.\ una restricción regulatoria o
de seguridad); de lo contrario, ignora compromisos intermedios mucho más
eficientes que sí captura el frente de Pareto.
